%% ma: L12-B13-0003
%% lop: 12
%% chuong: 4
%% bai: 13
%% dang: 12.13.1
%% loai: TN4
%% muc: NB
%% dap_an: "D"
%% nguon: "0. ĐỀ THAM KHẢO TN THPT 2025.pdf (D00), Phần I Câu 2"
%% kien_thuc: "Bài 13 (công thức thể tích khối tròn xoay quanh trục Ox)"
%% trang_thai: chinh-thuc
\begin{ex}
Cho hàm số $y=f(x)$ liên tục, nhận giá trị dương trên đoạn $[a;b]$. Xét hình phẳng $(H)$ giới hạn bởi đồ thị hàm số $y=f(x)$, trục hoành và hai đường thẳng $x=a$, $x=b$. Khối tròn xoay được tạo thành khi quay hình phẳng $(H)$ quanh trục $Ox$ có thể tích là
\choice{$V=\pi\displaystyle\int_a^b\big|f(x)\big|\,\mathrm dx$}{$V=\pi^2\displaystyle\int_a^bf(x)\,\mathrm dx$}{$V=\pi^2\displaystyle\int_a^b\big[f(x)\big]^2\,\mathrm dx$}{\True $V=\pi\displaystyle\int_a^b\big[f(x)\big]^2\,\mathrm dx$}
\loigiai{Theo công thức thể tích khối tròn xoay sinh bởi hình phẳng quay quanh $Ox$: $V=\pi\displaystyle\int_a^b\big[f(x)\big]^2\,\mathrm dx$.}
\end{ex}
